\section{Experiments and Quantitative Sweeps}
\label{sec:experiments}

To validate our framework with maximum empirical rigor, we execute three extensive experimental sweeps over 5 independent random seeds. We analyze: (1) sensitivity to the confidence threshold $\gamma_{\text{conf}}$, (2) stability under calibration sample scarcity $N$, and (3) resilience against batch heterogeneity collapse in production streams.

\subsection{Experimental Setup}
We conduct all experiments using our pre-calibrated \emph{Isolating Coordinate Sandbox}. There are $K=4$ experts, each specializing in one of the tasks: MNIST (Task 0), Fashion-MNIST (Task 1), CIFAR-10 (Task 2), and SVHN (Task 3). The pre-calibrated noise scales $\{\sigma_0=0.05, \sigma_1=0.05, \sigma_2=0.35, \sigma_3=1.25\}$ set the standalone performance ceilings for each expert.
The parametric routing networks are trained using the AdamW optimizer with a learning rate of $1\times 10^{-3}$, a batch size of 256, and $L_2$ weight decay ($0.1$) for 100 epochs. All experiments are averaged across 5 independent random seeds, and we report the mean and standard deviation.

We critically note that while our static baselines include only Uniform Merging, advanced static model-merging techniques such as Task Arithmetic, TIES-Merging, and DARE mathematically reduce exactly to Uniform Merging in our coordinate-isolated sandbox. Since different task experts reside in disjoint blocks of coordinate dimensions (non-overlapping weights), there are no conflicting parameter deltas across different experts. Consequently, sign-agreement checks (like in TIES-Merging) or random weight dropping (like in DARE) do not prune or modify any overlapping coordinate conflicts. After standard temperature or scaling normalization, these advanced static ensembling methods behave identically to a uniform or scaled-uniform weight blend, making Uniform Merging the mathematically exact representative of all static model-merging baselines in this environment.

\subsection{Main Quantitative Results}
Table \ref{tab:main_results} presents the performance of all baselines compared to our proposed CGHR under standard homogeneous batching ($B=256$) with calibration size $N=64$.

\begin{table*}[t]
\caption{Multi-task performance comparison on the Isolating Coordinate Sandbox under standard Homogeneous Batching ($B=256$, calibration size $N=64$). All statistics represent the mean and standard deviation computed across 5 independent random seeds.}
\label{tab:main_results}
\vskip 0.15in
\begin{center}
\begin{small}
\setlength{\tabcolsep}{3.5pt}
\begin{tabular}{lcccccc}
\toprule
Method & Trainable Params & MNIST (\%) & F-MNIST (\%) & CIFAR-10 (\%) & SVHN (\%) & Joint Mean (\%) \\
\midrule
Expert Ceiling & 0 & 100.00 & 100.00 & 88.60 & 26.40 & 78.75 \\
Uniform Merging & 0 & 95.20 & 95.00 & 45.00 & 17.20 & 63.10 $\pm$ 0.00 \\
\midrule
Linear Router (Unreg) & 772 & 99.80 & 100.00 & 54.80 & 18.80 & 68.74 $\pm$ 1.23 \\
Linear Router (Reg) & 772 & 99.80 & 100.00 & 54.80 & 18.80 & 68.66 $\pm$ 1.15 \\
VR-Router & 772 & 99.80 & 100.00 & 54.80 & 18.80 & 68.63 $\pm$ 1.16 \\
TSAR & 772 & 97.00 & 97.60 & 47.20 & 18.00 & 65.18 $\pm$ 0.15 \\
\midrule
PFSR & 0 & 100.00 & 100.00 & 88.00 & 18.40 & 76.60 $\pm$ 0.00 \\
CGHR (Ours) & 772 & 100.00 & 100.00 & 87.00 & 18.80 & \textbf{76.44} $\pm$ \textbf{0.09} \\
\bottomrule
\end{tabular}
\end{small}
\end{center}
\vskip -0.1in
\end{table*}

Our analysis reveals several critical empirical insights:
\begin{itemize}
    \item \textbf{Parametric Overfitting on Difficult Tasks}: Standard parametric models (Linear, VR-Router) achieve excellent results on simple tasks like MNIST and Fashion-MNIST but experience severe overfitting on the more difficult CIFAR-10 dataset ($54.80\%$ compared to the $88.60\%$ ceiling). This is because the linear decision boundary cannot generalize robustly given limited calibration data.
    \item \textbf{The Power of Parameter-Free Projections}: Pure PFSR achieves outstanding zero-shot performance ($76.60\%$), performing within $0.60\%$ of the CIFAR-10 ceiling without updating any weights. This demonstrates the high empirical quality of cosine similarity projections onto the frozen expert classification manifolds.
    \item \textbf{High Stability of CGHR}: Our proposed CGHR matches PFSR closely on MNIST, F-MNIST, and CIFAR-10, and outperforms it on SVHN ($18.80\%$ vs. $18.40\%$). Critically, CGHR exhibits extremely low variance across 5 independent seeds ($\pm 0.09$), delivering a highly stable and robust compromise.
    \item \textbf{Cascaded Routing Failure on Weak Experts}: Under SVHN (Task 3), the standalone expert ceiling is low ($26.40\%$) due to high noise ($\sigma_3 = 1.25$). Both PFSR ($18.40\%$) and CGHR ($18.80\%$) perform significantly below this ceiling. This represents a classic \emph{cascaded routing failure}: when an expert is highly noisy and weak, its incoming representations are highly corrupted, which degrades both the similarity projections and parametric gating logits. Consequently, the gateway router frequently misclassifies the active task (often routing SVHN inputs to cleaner, highly precise experts like CIFAR-10), ensuring downstream classification failure.
\end{itemize}

\subsection{Empirical Sweep 1: Confidence Gating Threshold Sensitivity}
We evaluate the sensitivity of CGHR to the gating threshold $\gamma_{\text{conf}} \in [0.0, 1.0]$. A threshold of $0.0$ corresponds to a pure parametric router (Linear Reg), while a threshold of $1.0$ corresponds to pure PFSR. We sweep this threshold across three different confidence metrics: Max Probability, Negative Entropy, and Margin.

\begin{figure}[htbp]
  \centering
  \includegraphics[width=0.9\linewidth]{fig1_confidence_sweep.png}
  \caption{Confidence gating threshold sweep ($\gamma_{\text{conf}} \in [0.0, 1.0]$) across Max Probability, Negative Entropy, and Margin confidence metrics under heterogeneous batching without MBH. An optimal peak performance envelope is observed around $\gamma_{\text{conf}} = 0.85$ for Max Probability.}
  \label{fig:conf_sweep}
\end{figure}

As shown in Figure \ref{fig:conf_sweep}:
\begin{itemize}
    \item There exists a robust \textbf{peak performance envelope} at intermediate thresholds between $0.75$ and $0.95$.
    \item Under Max Probability, the model achieves peak Joint Mean accuracy at $\gamma_{\text{conf}} = 0.85$. This confirms that confidence gating successfully routes low-confidence, ambiguous, or OOD samples to the robust PFSR manifold while preserving highly precise parametric decisions for high-confidence, in-distribution samples.
    \item Margin and Negative Entropy show similar peak envelopes, confirming that the gating benefit is robust across different confidence metrics.
\end{itemize}

\subsection{Empirical Sweep 2: Generalization Under Data Scarcity}
To test the transductive overfitting resilience of all routers, we sweep the calibration sample complexity $N$ across $\{16, 32, 64, 128, 256, 512\}$.

\begin{figure}[htbp]
  \centering
  \includegraphics[width=0.9\linewidth]{fig2_sample_complexity.png}
  \caption{Calibration sample complexity sweep ($N \in \{16, 32, 64, 128, 256, 512\}$) across 5 independent seeds. CGHR (Ours) exhibits flatline-like stability, leveraging the PFSR fallback to prevent transductive overfitting at small $N$ while absorbing parametric gains as $N$ increases.}
  \label{fig:sample_complexity}
\end{figure}

As plotted in Figure \ref{fig:sample_complexity}:
\begin{itemize}
    \item At extremely low data regimes ($N = 16$), unregularized linear routers overfit severely and display massive variance across seeds. Regularization ($L_2$ weight decay, VR, and TSAR) helps stabilize training but still results in substantial performance degradation compared to larger $N$.
    \item Pure PFSR remains perfectly flat because it is zero-shot and requires zero calibration data.
    \item Our proposed \textbf{CGHR (Ours)} demonstrates outstanding flatline-like stability across all sample complexities. At $N=16$, CGHR leverages the robust PFSR fallback to maintain high performance, and as $N$ grows, it seamlessly integrates the refined parametric linear updates to scale up and match or exceed the regularized baselines. We note that optimizing the gating threshold $\gamma_{\text{conf}}$ under extreme data scarcity like $N=16$ presents a unique \emph{Calibration Paradox} since setting aside a separate validation partition is highly impractical, while tuning directly on the training pool risks severe overfitting. To resolve this paradox for low-resource practitioners, we elevate three highly practical, training-free, and data-efficient calibration strategies detailed in Appendix~\ref{app:calibration_paradox}:
    \begin{itemize}
        \item \textbf{High-Dimensional Random Projection Prior (Zero-Data Calibration)}: Leveraging high-dimensional random projection theory, we prove that the expected maximum cosine similarity of a purely random, out-of-distribution (OOD) vector with $K$ experts in a $D$-dimensional space scales as $\mathcal{O}(\sqrt{2\log K / D})$. By projecting Gaussian random vectors onto the frozen expert registry offline, we can construct a robust, data-free threshold $\gamma_{\text{conf}}^{\text{prior}} = \mu_{\text{random}} + z_{1-\alpha} \sigma_{\text{random}}$ (where $z_{1-\alpha}$ is the standard normal critical value) without consuming a single calibration sample.
        \item \textbf{Leave-One-Out Cross-Validation (LOO-CV)}: Since our parametric linear router is extremely lightweight and trains in $<5$ ms, we can execute LOO-CV over the small calibration set in under $150$ ms. This provides an unbiased estimate of the generalization gating confidence without sacrificing a single sample from the training pool.
        \item \textbf{Self-Calibrating Unsupervised Stream Gating}: Instead of a static threshold, we can calibrate $\gamma_{\text{conf}}(t) = \gamma_{\text{base}} + \eta \bar{H}_t$ dynamically at runtime using the rolling average entropy $\bar{H}_t$ of the incoming deployment stream, automatically shifting routing towards the robust parameter-free fallback as stream volatility increases.
    \end{itemize}
\end{itemize}

\subsection{Empirical Sweep 3: Robustness to Heterogeneity Collapse}
To simulate realistic deployment conditions, we subject the dynamic routers to mixed-task heterogeneous streams and sweep the batch size $B$ across $\{1, 8, 32, 128, 512\}$.

\begin{figure}[htbp]
  \centering
  \includegraphics[width=0.9\linewidth]{fig3_stream_audit.png}
  \caption{Deployment stream audit sweeping batch size ($B \in \{1, 8, 32, 128, 512\}$) under mixed-task heterogeneous streams. Without MBH, all standard routers experience severe heterogeneity collapse. Integrating MBH completely prevents collapse, keeping performance flat and robust.}
  \label{fig:stream_audit}
\end{figure}

As shown in Figure \ref{fig:stream_audit}:
\begin{itemize}
    \item Without MBH, all standard routing architectures experience severe \emph{heterogeneity collapse}. As the batch size increases, the joint mean accuracy of standard parametric and PFSR routers degrades rapidly, eventually collapsing to the performance of Uniform Merging ($63.10\%$) at $B=512$. This occurs because mixed-task samples smooth out the routing activations over the entire batch, preventing task isolation.
    \item In stark contrast, when combining routers with our proposed \textbf{Micro-Batch Homogenization (MBH)}, both \textbf{PFSR + MBH} and \textbf{CGHR + MBH} exhibit perfectly flat, robust, collapse-free performance curves across all batch sizes. By dynamically grouping and isolating homogeneous sub-batches, MBH completely protects the expert fusion boundaries and maintains high performance under mixed deployment streams.
\end{itemize}

\subsection{Structural Asymmetry and Sandbox Limitations}
\label{sec:structural_asymmetry}

A critical, rigorous auditing of our synthetic sandbox design reveals an important structural asymmetry between Pathway A (Parametric Gating) and Pathway B (PFSR). This asymmetry, which we examine with scientific transparency, centers on the routing inputs:
\begin{itemize}
    \item \textbf{Parametric Input Space}: The trained parametric linear router (Pathway A) takes the global $D$-dimensional feature vector $z_b \in \mathbb{R}^{192}$ as input (Equation 1). Consequently, it must learn to map features to task outputs while filtering out high-variance Gaussian noise in the other $K-1$ non-active block coordinates. Under scarce calibration data ($N \le 64$), this high-dimensional noise leads to transductive overfitting, which artificially degrades its performance on complex datasets like CIFAR-10 ($54.80\%$).
    \item \textbf{Non-Parametric Input Space}: The parameter-free subspace router (Pathway B) is formulated to project block-specific representations $z_{k, b} \in \mathbb{R}^{48}$ onto Expert $k$'s classification manifold (Equation 5). This means PFSR is equipped with privileged architectural knowledge regarding the coordinate boundaries of each task.
\end{itemize}

This structural design is a deliberate choice made to model a scenario where a non-parametric fallback is aided by known topological expert-feature mappings (e.g., when experts are assigned to localized network modules or specific prompt tokens), whereas the trained router operates with general global views. Specifically, we note that the gating router's fundamental objective is to identify the active task coordinate from the global representation $z_b \in \mathbb{R}^{192}$. If the parametric router were pre-provided with the active block coordinate, the routing task itself would be trivially solved, as there would be no coordinate selection problem left to learn. Therefore, the parametric router \emph{must} receive the global feature vector $z_b$ as input to fulfill its function as an adaptive router, whereas the parameter-free router (PFSR) relies on local projection coordinates precisely because it operates as a localized subspace mapping heuristic. 

\subsection{Advanced Analyses and Theoretical Extensions}
\label{sec:advanced_analyses}

To preserve a highly focused and concise main body, several key mathematical formulations, comprehensive system audits, and robust error propagation sweeps are detailed in the Appendices:
\begin{itemize}
    \item \textbf{The UNC-PFSR Equivalence Theorem}: We prove a fundamental equivalence showing that under Unit-Norm Calibration (UNC), Global unpartitioned and Local block-sliced similarities are mathematically and empirically identical, resolving the coordinate-asymmetry of our sandbox. We also evaluate a lightweight test-time calibration protocol called \emph{Inference-Time block-wise Unit-Norm Calibration (IT-UNC)} that completely recovers Global PFSR accuracy under coordinate noise, rising from $30.00\%$ to the optimal $74.20\%$ ceiling (see Appendix~\ref{app:unc_equivalence}).
    \item \textbf{Cascaded Error Propagation and Mitigations}: We systematically corrupt the gateway router's accuracy using simulated error rates $P_{\text{error}} \in [0.0, 0.75]$. We analyze the resulting cascaded downstream accuracy collapses and empirically validate two highly practical design mitigations: \emph{Soft-Confidence Fallback Homogenization} (which blends routing towards a safe uniform prior under high uncertainty) and \emph{Hierarchical MBH} (which clusters experts into deterministically stable weight-space groups), demonstrating that they successfully buffer the serving engine from routing errors (see Appendix~\ref{app:error_propagation_detailed}).
    \item \textbf{Systems Latency and Hardware Realism}: We address CPU-bound Python-loop benchmarking simulation artifacts, modeling a vectorized parallel GPU environment exhibiting a typical $2\times$ to $4\times$ parallel occupancy latency penalty, and we outline a custom hardware-native Triton Segmented-BGEMM kernel design to bypass sequential serving overheads (see Appendix~\ref{app:systems_analysis}).
\end{itemize}